\section{本章小结}
\label{sec:sobolev-summary}

本章开头的问题是：怎样的度量既能感知振荡的快慢，又允许折点这样的适度不光滑。回答是同时控制函数与它的弱导数。$\sin(2\pi nx)$说明$L^2$大小相同的函数可以有越来越大的导数，常数说明只看导数会丢掉函数的整体水平；$|x-1/2|$有$L^2$弱导数而属于$H^1(0,1)$，阶跃函数的跳跃产生点值作用，没有弱导数，几乎处处导数也不能代替弱导数。Sobolev范数把函数及直到$k$阶的弱导数的$L^p$大小合在一起，使每个求导映射$W^{k,p}\to L^p$连续。完备性的关键是测试恒等式在$L^p$极限下保持，因而各分量的极限正是极限函数的弱导数；$p=2$时再得到Hilbert结构。

空间中的操作各有自己的条件。对任意开集与有限指数，局部磨光、截断与局部有限的单位分解拼接出内部光滑的逼近函数，但它们不一定紧支撑，也不一定光滑到边界；$p=\infty$时这种逼近失败。跨边界延拓在平直边界上用反射实现，在有界Lipschitz区域上通过局部拉直拼接成有界线性算子，并由此得到光滑到边界的逼近。零延拓需要边界值为零，高阶延拓需要各阶法向导数在边界两侧相接。

边界约束先由紧支撑光滑函数的闭包$W_0^{k,p}$表达，它是闭线性子空间，不涉及任何边界值。一维时，绝对连续代表的端点值受$H^1$范数控制，因而$H_0^1$中的函数端点为零，常数$1$不在其中。高维时，迹映射先对光滑到边界的函数限制到边界，再借逼近延伸为$W^{1,p}\to L^p(\partial\Omega)$的有界算子；它不是逐点边界值，一阶零迹恰好刻画$W_0^{1,p}$，高阶零边界还要求法向导数为零。排除常数方向以后，Poincaré不等式让导数单独给出等价范数：零边界情形只需区域在一个方向上有界，零均值情形还需要区域连通。

连续嵌入回答单个函数能被控制到什么程度，答案由$p$与$d$的比较决定。$p<d$时得到更高的可积性$L^{p^*}$，幂函数说明$p^*$不能提高，且函数可以无界；$p=d\geq2$时对每个有限指数都有估计，但$\ln\ln(e/|x|)$说明没有统一的有界性；$p>d$时有$1-d/p$阶Hölder连续代表，$p=\infty$时有Lipschitz代表。一维的$W^{1,1}$已经有有界的绝对连续代表，高维的临界结论不能照搬到一维。

紧嵌入回答有界函数列的行为。在有界Lipschitz区域上，延拓、平移差的统一控制与有限维近似给出$L^p$收敛子列，再由插值得到$q<p^*$时的$L^q$收敛子列；这种收敛不包括导数，因而不是Sobolev范数下的收敛。临界指数处，把固定形状压向一点的集中序列保持$L^{p^*}$范数不变而几乎处处趋零，所以没有收敛子列，严格指数条件不能去掉。转写为函数类，紧嵌入说明Sobolev球在次临界$L^q$距离下对每个精度都有有限覆盖，但只给出数量有限，不给出增长阶；临界处的连续嵌入与弱紧性都不能代替这一结论。

本章的整数阶空间还不能描述阶跃函数那样“差一点就有导数”的变化，也没有区分变化在不同尺度上怎样分布。第~\ref{chap:besov-spaces}章用平移差与小波系数补上这两点，并在那里算出覆盖数的增长阶；具有连续点值读取的Hilbert函数空间则是第~\ref{chap:rkhs}章的对象。
